\section{Homogeneous transport, anisotropic dispersion, and nonlocal attraction}
\label{sec:r11-transport}

The transport argument extends to general homogeneous symbols without a
sign assumption on the localized principal commutator. Its key is an
exact cancellation on a spatially supported core. We obtain a common
time-dependent result for local critical powers and Hartree attraction.
The quantity localized in the general theorem is the full homogeneous
kinetic density; the stronger gradient conclusion for isotropic quartic
dispersion is considered separately at the end.

\subsection{Equations and solution class}
Fix $m>1$, $N\ge1$, and $s=m/2$. Let $a$ be real, even, homogeneous
of degree $m$, $C^{N+2}$ away from zero, and positive elliptic:
\[
 c|\xi|^m\le a(\xi)\le C|\xi|^m.
\]
Set
\[
 A(u)=\int a(\xi)|\widehat u(\xi)|^2\,d\xi,
 \quad M(u)=\|u\|_2^2,
 \quad P(u)=\int\xi|\widehat u(\xi)|^2\,d\xi.
\]
Momentum is finite since $m>1$. We treat either
\begin{align}
 iu_t&=a(D)u-|u|^{p-2}u,
 &p&=2+2m/N,
 &E(u)&=\tfrac12A(u)-\tfrac1p\|u\|_p^p,
 \label{r11tr:local}
\end{align}
or, for $1<m<N$,
\begin{align}
 iu_t&=a(D)u-(K*|u|^2)u,
 &E(u)&=\tfrac12A(u)-\tfrac14H(u),
 \label{r11tr:hartree}\\
 H(u)&=\iint K(x-y)|u(x)|^2|u(y)|^2\,dx\,dy.
 \notag
\end{align}
Here $K$ is real, even, $C^1$ off zero, homogeneous of degree $-m$, and
\begin{equation}\label{r11tr:kernel}
 |K(x)|+|x||\nabla K(x)|\le C_K|x|^{-m}.
\end{equation}
Positive type is unnecessary for the transport theorem: all nonlinear
error estimates use the absolute Riesz interaction. This differs from
the positivity needed for the phase-kernel argument.

Throughout this section, $u\in C([0,\infty);H^s)$ is a strong energy
solution with conserved $M,P,E$, satisfying
\begin{equation}\label{r11tr:bound}
 \sup_{t\ge0}\|u(t)\|_{H^s}\le B<\infty.
\end{equation}
We assume the solution is obtained, on compact time intervals, as a
strong energy-space limit of sufficiently smooth solutions, so that
the integrated localized virial identity follows from the estimates
below. No local well-posedness theorem for arbitrary anisotropic
symbols is asserted here.

Choose $\chi\in C_c^\infty$ with $0\le\chi\le1$, equal to one on
$|x|\le1$ and zero on $|x|\ge2$, and put $\eta=1-\chi$.
For a center $c(t)$ define
\begin{align}
 T_R^c(t)&=\bigl\||D|^s[\eta((x-c(t))/R)u(t)]\bigr\|_2^2,
 \label{r11tr:cut-tail}\\
 D_R^c(t)&=\int_{|x-c(t)|>R}\bigl||D|^s u(t,x)\bigr|^2\,dx.
 \label{r11tr:literal-tail}
\end{align}
These are distinct quantities, but the cutoff commutator below makes
their localization properties equivalent on a bounded energy ball.

\subsection{The kinetic commutator is exact on the core}
Choose a real smooth compactly supported vector field $f$ equal to $x$
on $|x|\le4$. Set
\[
 f_R(x)=Rf(x/R),\qquad
 L_R=f_R\cdot\nabla+\tfrac12\operatorname{div}f_R,
 \qquad B_R=-iL_R,
 \qquad C_R=[a(D),L_R].
\]
Thus $L_R$ is skew-adjoint, while $B_R$ and $C_R$ are symmetric forms.
Write
\[
 Q_R(w)^2=\||D|^s w\|_2^2+R^{-m}\|w\|_2^2.
\]

\begin{lemma}[Uniform form bound and core cancellation]
\label{lem:r11tr-core}
For all $R>0$ and $z,w\in H^s$,
\begin{equation}\label{r11tr:form}
 |\langle z,C_Rw\rangle|+
 |\langle z,ma(D)w\rangle|\le C Q_R(z)Q_R(w).
\end{equation}
If $z$ is supported in $|x|\le2R$, then
\begin{equation}\label{r11tr:exact-core}
 \langle z,[C_R-ma(D)]z\rangle=0.
\end{equation}
Consequently, on a bounded $H^s$ ball and for $R\ge1$,
\begin{equation}\label{r11tr:kinetic-tail}
 \left|\langle u,C_Ru\rangle-mA(u)\right|
 \le C_B\left(\||D|^s(\eta_Ru)\|_2+R^{-s}\right).
\end{equation}
All statements also hold about any fixed center.
\end{lemma}
\begin{proof}
At $R=1$ the Fourier kernel of $C_1$, up to a fixed normalization and
the factor $i$, is
\[
 \tfrac12[a(\xi)-a(\zeta)](\xi+\zeta)\cdot
                   \widehat f(\xi-\zeta).
\]
The homogeneous derivative bound gives
\[
 |a(\xi)-a(\zeta)|\,|\xi+\zeta|
 \le C|\xi-\zeta|(|\xi|+|\zeta|)^m.
\]
After division by $\langle\xi\rangle^s\langle\zeta\rangle^s$,
the right side is at most
$C|\xi-\zeta|\langle\xi-\zeta\rangle^s$. Its product with
$|\widehat f(\xi-\zeta)|$ is integrable. Schur's test therefore
bounds $C_1:H^s\to H^{-s}$. Under the unitary rescaling
$v(y)=R^{N/2}w(Ry)$, $C_R$ becomes $R^{-m}C_1$.
This proves \eqref{r11tr:form}; the term $ma(D)$ is immediate.

For smooth compactly supported $z$ in the stated ball,
$L_Rz=(x\cdot\nabla+N/2)z$. Homogeneity and differentiation of
$A(\lambda^{N/2}z(\lambda\cdot))=\lambda^mA(z)$ give
\[
 2\operatorname{Re}\langle a(D)z,L_Rz\rangle=mA(z),
\]
which is \eqref{r11tr:exact-core}. Approximation within a slightly
larger ball and \eqref{r11tr:form} extend the equality to supported
$H^s$ functions. This is equality of quadratic forms, not a claim
that the defect operator annihilates $z$ pointwise.

Let $z=\chi_Ru$, $w=\eta_Ru$ and $\mathcal D_R=C_R-ma(D)$.
The core equality gives
\[
 \langle u,\mathcal D_Ru\rangle
 =\langle w,\mathcal D_Rw\rangle
       +2\operatorname{Re}\langle z,\mathcal D_Rw\rangle.
\]
The smooth-cutoff product estimate
\[
 \||D|^s(\chi_Ru)\|_2
 \le C\bigl(\||D|^su\|_2+R^{-s}\|u\|_2\bigr)
\]
and \eqref{r11tr:form} now prove \eqref{r11tr:kinetic-tail}.
The product estimate itself follows at scale one from the convolution
inequality $|\xi|^s\le C_s(|\zeta|^s+|\xi-\zeta|^s)$ and Young's
inequality, followed by rescaling.
\end{proof}

No favorable sign of the principal symbol
$\nabla a(\xi)\cdot Df_R(x)\xi$ is used. In particular, the proof
does not infer such a sign from ellipticity or convexity of $a$.

\begin{lemma}[Cutoff energy and the literal kinetic tail]
\label{lem:r11tr-tail-equivalence}
Put $\delta=\min(s,1)$. On the ball \eqref{r11tr:bound},
\begin{equation}\label{r11tr:tail-equivalence}
 (T_R^c)^{1/2}\le(D_R^c)^{1/2}+C_BR^{-\delta},\qquad
 (D_{2R}^c)^{1/2}\le(T_R^c)^{1/2}+C_BR^{-\delta}.
\end{equation}
\end{lemma}
\begin{proof}
For $0<s<1$, the Fourier inequality
$||\xi|^s-|\zeta|^s|\le|\xi-\zeta|^s$ and Young's inequality give
\[
 \|[|D|^s,\eta_R]u\|_2\le CR^{-s}\|u\|_2.
\]
For $s\ge1$, use
\[
 ||\xi|^s-|\zeta|^s|
 \le C_s|\xi-\zeta|
             (|\zeta|^{s-1}+|\xi-\zeta|^{s-1})
\]
to obtain
\[
 \|[|D|^s,\eta_R]u\|_2
 \le C\left(R^{-1}\||D|^{s-1}u\|_2+R^{-s}\|u\|_2\right).
\]
Only the compact part $-\chi_R$ contributes to the commutator.
The uniform bound is $C_BR^{-\delta}$. Since $\eta_R=0$ inside
the radius-$R$ ball and $\eta_R=1$ outside the radius-$2R$ ball,
the triangle inequality proves both assertions.
\end{proof}

\subsection{The common time-dependent estimate}
Let $c$ be Lipschitz with $\|c'\|_\infty\le L$ and set
\begin{equation}\label{r11tr:virial}
 J_R(t)=2\langle u(t),B_R^{c(t)}u(t)\rangle
 =2\operatorname{Im}\int f_R(x-c(t))\cdot\overline u\nabla u.
\end{equation}
For $s<1$, the last expression denotes the symmetric $H^{1/2}$
quadratic form, defined by density.

\begin{lemma}[Two-sided localized virial estimate]
\label{lem:r11tr-virial}
Set
\[
 \gamma=\min\left(\tfrac12,\tfrac1m\right),\qquad
 \beta=\min\left(\tfrac m2,1\right),\qquad
 \widetilde\beta=\min\left(\tfrac m2,\tfrac2m\right).
\]
There are constants $C_0,C_1$ depending on the fixed equation,
$M,B,L$ and cutoffs such that, for $R\ge1$,
\begin{align}
 |J_R(t)|&\le C_0R,\label{r11tr:endpoint}\\
 |J_R'-4mE+2c'\cdot P|
 &\le C_1\left((T_R^c)^\gamma+R^{-\beta}\right),
 \label{r11tr:virial-cut}\\
 |J_R'-4mE+2c'\cdot P|
 &\le C_1\left((D_R^c)^\gamma+R^{-\widetilde\beta}\right).
 \label{r11tr:virial-literal}
\end{align}
The derivative assertions are understood as integrated identities with
an error bounded by the integral of the displayed right side.
\end{lemma}
\begin{proof}
The scale-one first-order Fourier-Schur form estimate and rescaling give
\[
 |\langle u,B_Ru\rangle|
 \le C\bigl(R\||D|^{1/2}u\|_2^2+M\bigr).
\]
Since $s>1/2$, this proves \eqref{r11tr:endpoint}. The kinetic
contribution to $J_R'$ is $2\langle u,C_Ru\rangle$.
Lemma~\ref{lem:r11tr-core} controls its error from $2mA(u)$.

For the local equation, direct integration by parts gives the nonlinear
term
\[
 -\frac{2(p-2)}p\int\operatorname{div}f_R(x-c)|u|^p.
\]
Its difference from $-4m\|u\|_p^p/p$ is bounded by
$C\int_{|x-c|>4R}|u|^p$. Critical interpolation gives
\begin{equation}\label{r11tr:local-tail}
 \int_{|x-c|>4R}|u|^p
 \le\|\eta_Ru\|_p^p\le C M^{m/N}T_R^c.
\end{equation}

For the Hartree equation, write $\rho=|u|^2$. Symmetrization gives
the nonlinear term
\[
 \iint\rho(x)\rho(y)
  [f_R(x-c)-f_R(y-c)]\cdot\nabla K(x-y)\,dx\,dy.
\]
Its difference from $-mH(u)$ vanishes if both points lie in
$B(c,4R)$. The Lipschitz bound on $f_R$ and
\eqref{r11tr:kernel} dominate the error by the absolute Riesz
interaction with at least one point outside that ball. For
$q=4N/(2N-m)$, Hardy--Littlewood--Sobolev and interpolation yield
\begin{align}
 |\text{Hartree error}|
 &\le C\|\eta_Ru\|_q^2\|u\|_q^2
 \le C_{M,B}(T_R^c)^{1/2},\label{r11tr:hartree-tail}\\
 \|w\|_q^2&\le C\|w\|_2\||D|^sw\|_2.\notag
\end{align}
Symmetrization supplies one factor $|x-y|$ at the diagonal, so HLS
also justifies passage from a regularized kernel.

Motion of the center adds
\[
 -2c'(t)\cdot\int Df_R(x-c)\operatorname{Im}(\overline u\nabla u).
\]
The difference from $-2c'\cdot P$ uses $g_R=Df_R-I$, which
vanishes on $B(c,4R)$. Its current form on $u$ is exactly its
current form on $w=\eta_Ru$: the cutoff is one on the support
of $g_R$, and its derivative contributes only a real multiple of
$|u|^2$. The rescaled first-order form estimate gives
\begin{align*}
 \left|\operatorname{Im}\int g_R\overline w\nabla w\right|
 &\le C\bigl(\||D|^{1/2}w\|_2^2+R^{-1}\|w\|_2^2\bigr)\\
 &\le C_M(T_R^c)^{1/m}+C_MR^{-1}.
\end{align*}
Here $g_R$ is a constant matrix plus a rescaled smooth compactly
supported matrix; its constant part is the ordinary momentum form.
The last inequality is interpolation between orders zero and $s$.
All current identities extend by density in $H^{1/2}$.

The principal terms are $4mE-2c'\cdot P$ in both equations.
Since $T_R^c$ is uniformly bounded, reducing the powers
$1/2,1,1/m$ to $\gamma$ proves \eqref{r11tr:virial-cut}.
Lemma~\ref{lem:r11tr-tail-equivalence} gives
\[
 (T_R^c)^\gamma\le C(D_R^c)^\gamma+C_BR^{-2\delta\gamma}.
\]
Since $2\delta\gamma=\widetilde\beta\le\beta$, this proves
\eqref{r11tr:virial-literal}. Every form and nonlinear expression
used above is continuous on bounded energy-space sets. Approximation
therefore proves the integrated statements for the solution class
specified above.
\end{proof}

\subsection{Persistent exterior kinetic density and exact averaged drift}

\begin{theorem}[Transport at a distance proportional to the time window]
\label{thm:r11tr-ballistic}
Assume $E=-e<0$ and $P=0$. For every speed bound $L$, there are
$\nu,d>0$, depending only on the fixed equation, $M,B,e,L$ and
cutoffs, such that every center $c$ with speed at most $L$ satisfies
\begin{equation}\label{r11tr:ballistic}
 \liminf_{T\to\infty}\inf_{S\ge0}\frac1T
  \int_S^{S+T}\int_{|x-c(t)|>\nu T}
                \bigl||D|^{m/2}u(t,x)\bigr|^2\,dx\,dt
 \ge d.
\end{equation}
\end{theorem}
\begin{proof}
Integrate the upper half of \eqref{r11tr:virial-literal} over
$[S,S+T]$ and use \eqref{r11tr:endpoint}:
\[
 C_1\frac1T\int_S^{S+T}(D_R^c)^\gamma\,dt
 \ge4me-C_1R^{-\widetilde\beta}-2C_0R/T.
\]
Choose $\nu>0$ with $2C_0\nu\le me$ and put $R=\nu T$,
held fixed during this time integration. For large $T$, the right
side is at least $2me$. Jensen's inequality, using $\gamma\le1$,
gives the asserted lower bound with $d=(2me/C_1)^{1/\gamma}$.
The bounds are uniform in $S$ and in all centers of the given speed.
\end{proof}

\begin{remark}
The conclusion concerns the nonlocal homogeneous kinetic density when
$m$ is not an even integer. Its exterior integral is uniformly bounded,
so the lower average forces a positive fraction of times in each long
window to have nonzero exterior kinetic density. This is not a
finite-speed theorem or a finite radiation-budget estimate, and does
not exclude indefinite replenishment of exterior activity.
\end{remark}

\begin{theorem}[Exact drift under averaged kinetic localization]
\label{thm:r11tr-drift}
For either equation and a Lipschitz center $c$, suppose
\begin{equation}\label{r11tr:average-tight}
 \lim_{R\to\infty}\limsup_{T\to\infty}\frac1T
                  \int_0^T D_R^c(t)\,dt=0.
\end{equation}
Then the limit exists and
\begin{equation}\label{r11tr:drift}
 \lim_{T\to\infty}\frac{P\cdot[c(T)-c(0)]}{T}=2mE.
\end{equation}
In particular $E<0$ and $P=0$ exclude
\eqref{r11tr:average-tight}.
\end{theorem}
\begin{proof}
Integrate the two-sided estimate \eqref{r11tr:virial-literal},
divide by $T$, and use \eqref{r11tr:endpoint} and Jensen.
For fixed $R$, the limiting discrepancy from \eqref{r11tr:drift}
is at most
\[
 C\left[R^{-\widetilde\beta}+
       \left(\limsup_{T\to\infty}\frac1T
                    \int_0^TD_R^c(t)\,dt\right)^\gamma\right].
\]
Let $R\to\infty$.
\end{proof}

The drift theorem holds at every energy sign. For a traveling wave,
it reproduces $v\cdot P=2mE$. Its content for a general localized
trajectory does not require a rigid profile or periodicity.

\subsection{Uniform localization with an arbitrary center}

\begin{theorem}[No regularity assumption on the localization center]
\label{thm:r11tr-center}
Suppose $E<0$ and an arbitrary map $b:[0,\infty)\to\mathbb R^N$
satisfies
\begin{equation}\label{r11tr:uniform-tight}
 \sup_{t\ge0}D_R^b(t)\longrightarrow0\qquad(R\to\infty).
\end{equation}
Then $b$ remains within a bounded distance of a Lipschitz center,
and
\[
 \lim_{T\to\infty}\frac{P\cdot[b(T)-b(0)]}{T}=2mE.
\]
Consequently no solution obeying \eqref{r11tr:bound},
$E<0$, $P=0$, and \eqref{r11tr:uniform-tight} exists.
\end{theorem}
\begin{proof}
Lemma~\ref{lem:r11tr-tail-equivalence} first gives uniform
localization of $T_R^b$. In the local case take $z=p$; negative
energy gives $\|u(t)\|_p^p\ge p|E|$. In the Hartree case take
$z=q=4N/(2N-m)$. Since
$H(u)=2A(u)+4|E|\ge4|E|$ and
$|H(u)|\le C\|u\|_q^4$, the $L^q$ norm also has a uniform
positive lower bound. Estimates \eqref{r11tr:local-tail} and
\eqref{r11tr:hartree-tail}, applied at smaller cutoff radii,
show that the $L^z$ norm outside $B(b(t),2R)$ tends to zero
uniformly in time.

The equation gives $\|u_t\|_{H^{-s}}\le C_B$. For the local term,
this follows by H\"older and $H^s\hookrightarrow L^p$; for Hartree,
use
\[
 \left|\int(K*|u|^2)u\overline h\right|
 \le C\|u\|_q^3\|h\|_q\le C_B\|h\|_{H^s}.
\]
Choose $\sigma$ strictly between $N(1/2-1/z)$ and $s$.
Sobolev embedding and interpolation yield
\begin{equation}\label{r11tr:time-continuity}
 \|u(t)-u(t')\|_{L^z}
 \le C_B|t-t'|^{(s-\sigma)/(2s)}.
\end{equation}
Fix $R_0$ so that the exterior $L^z$ norm is much smaller than
the positive core norm, and then $\tau>0$ so that
\eqref{r11tr:time-continuity} is smaller still for
$|t-t'|\le\tau$. If $|b(t)-b(t')|>4R_0$, the two balls of
radius $2R_0$ are disjoint. Restriction to the core ball at $t'$
then contradicts \eqref{r11tr:time-continuity}. Hence
$|b(t)-b(t')|\le4R_0$ whenever $|t-t'|\le\tau$.

Linear interpolation through $b(n\tau)$ produces a center $c$
with $\|c'\|_\infty\le4R_0/\tau$ and
$\sup_t|c(t)-b(t)|\le8R_0$. A bounded displacement preserves
literal kinetic-tail localization by inclusion of balls. Apply
Theorem~\ref{thm:r11tr-drift}; the bounded displacement also
does not affect the drift limit.
\end{proof}

The localization assumption allows a diffuse fixed-mass tail whose
homogeneous kinetic norm tends to zero. Thus it is weaker than
precompactness in $H^s$ modulo translations. No uniform control of
high-frequency tails is needed.

\subsection{A stronger exterior-gradient theorem for quartic Hartree}

For $a(\xi)=|\xi|^4$, $N>4$, and a kernel satisfying
\eqref{r11tr:kernel} with $m=4$, put
\[
 G_R^c(t)=\int_{|x-c(t)|>R}|\nabla u(t,x)|^2\,dx.
\]
Use the usual radial cutoff $f_R=\nabla\phi_R$ with
$D^2\phi_R\le I$, equal to $x$ on the core. The biharmonic
kinetic term is bounded above by
$8\|\Delta u\|_2^2+O_B(R^{-2})$. In the Hartree estimate,
$q=2N/(N-2)$ and homogeneous Sobolev give
\[
 |\text{Hartree error}|
 \le C\|\eta_Ru\|_q^2\|u\|_q^2
 \le C_B\|\nabla(\eta_Ru)\|_2^2
 \le C_B(G_R^c+R^{-2}).
\]
The center-current error is at most
$C\sqrt M L(G_R^c)^{1/2}$. Consequently
\begin{equation}\label{r11tr:quartic-gradient}
 J_R'\le16E-2c'\cdot P+C\bigl((G_R^c)^{1/2}+R^{-2}\bigr).
\end{equation}

\begin{corollary}[Quartic nonlocal exterior-gradient persistence]
For the quartic Hartree equation just specified, a forward-global
uniformly $H^2$-bounded solution with $E<0$ and $P=0$ satisfies
\[
 \liminf_{T\to\infty}\inf_{S\ge0}\frac1T
 \int_S^{S+T}G_{\nu T}^c(t)\,dt\ge d>0
\]
for every bounded-speed center, with constants depending on its speed
bound. Uniform gradient localization about an arbitrary center is
impossible under those assumptions.
\end{corollary}
\begin{proof}
The time-window integration of \eqref{r11tr:quartic-gradient}
is identical to the proof of Theorem~\ref{thm:r11tr-ballistic}.
For an arbitrary uniformly localizing center, the $L^q$ lower bound,
the Sobolev estimate on the exterior gradient, and the temporal
continuity argument in Theorem~\ref{thm:r11tr-center} produce a
bounded-speed center. Fixed-radius integration then contradicts
$E<0,P=0$.
\end{proof}

This gradient conclusion uses the isotropic fourth-order kinetic sign;
it is not claimed for arbitrary anisotropic $a$. The general two-sided
theorem instead gives exact drift at the full kinetic-localization
level.

\subsection{Relation to earlier virial arguments}
Localized virial and compact-orbit rigidity are established methods.
Boulenger--Himmelsbach--Lenzmann, \emph{Blowup for fractional NLS},
arXiv:1509.08845, develop the fractional localized virial framework.
Dinh, arXiv:1808.07368, Proposition~1.5, already proves finite-time
blowup or an unbounded kinetic subsequence for nonradial mass-critical
isotropic dispersion of degree $1<m<2$ with negative energy, without
a localization hypothesis. That specialization is therefore not a
new open-problem resolution here. Cho--Hwang--Kwon--Lee,
arXiv:1208.2302, give earlier negative-energy Hartree blowup results
under their stated assumptions.

The extension established in this section is the common anisotropic
all-degree transport argument, including nonlocal kernels, the exact
averaged drift law and its kinetic-localization form. No priority claim
is based solely on this comparison. Nor do these results exclude a
bounded global orbit that indefinitely sustains the exterior activity
required by Theorem~\ref{thm:r11tr-ballistic}.
